\documentclass[11pt,a4paper]{article}

% ── Packages ──────────────────────────────────────────────────────────────────
\usepackage[utf8]{inputenc}
\usepackage[T1]{fontenc}
\usepackage[margin=2cm]{geometry}
\usepackage{enumitem}
\usepackage{booktabs}
\usepackage{array}
\usepackage{graphicx}
\usepackage{xcolor}
\usepackage{titlesec}
\usepackage{hyperref}
\usepackage{amsmath}
\usepackage{siunitx}
\usepackage{tikz}
\usepackage{fancybox}
\usepackage{tcolorbox}
\usetikzlibrary{shapes.geometric, arrows.meta, positioning, fit, calc}

\hypersetup{colorlinks=true, linkcolor=blue!60!black, urlcolor=blue!60!black}

\definecolor{guibg}{RGB}{240,248,255}
\definecolor{scriptbg}{RGB}{245,245,235}
\definecolor{warnbg}{RGB}{255,245,238}
\newtcolorbox{guistep}[1]{colback=guibg, colframe=blue!40!black, fonttitle=\bfseries, title=#1, boxrule=0.5pt, arc=2pt}
\newtcolorbox{scriptbox}[1]{colback=scriptbg, colframe=green!40!black, fonttitle=\bfseries, title=#1, boxrule=0.5pt, arc=2pt}
\newtcolorbox{warnbox}{colback=warnbg, colframe=red!50!black, fonttitle=\bfseries, title={Warning}, boxrule=0.5pt, arc=2pt}

\titleformat{\section}{\Large\bfseries}{Phase \thesection:~}{0em}{}
\titleformat{\subsection}{\large\bfseries}{Step \thesubsection~—~}{0em}{}
\titleformat{\subsubsection}{\normalsize\bfseries\itshape}{\thesubsubsection~}{0em}{}

\setlist[enumerate]{nosep, leftmargin=2em}
\setlist[itemize]{nosep, leftmargin=1.5em}

% ── Document ──────────────────────────────────────────────────────────────────
\begin{document}

\begin{center}
    {\LARGE\bfseries Plant Model — Detailed GUI Build Guide}\\[6pt]
    {\large BMS Engineering Assistant — Step-by-Step Simulink Instructions}\\[4pt]
    {\normalsize BMW i3 96s1p NMC Pack \quad|\quad MATLAB R2024a+ / Simulink}\\[10pt]
    \today
\end{center}

\vspace{4pt}
\noindent\rule{\textwidth}{0.8pt}

\tableofcontents
\newpage

%###############################################################################
%  PREREQUISITES
%###############################################################################
\section*{Prerequisites}
\addcontentsline{toc}{section}{Prerequisites}

\begin{guistep}{Before You Start}
\begin{enumerate}
    \item Open MATLAB. Set current folder to \texttt{plant\_model/}.
    \item In the MATLAB Command Window run:
    \begin{verbatim}
    run('scripts/export_luts.m')
    run('scripts/export_drivecycles.m')
    \end{verbatim}
    This loads \texttt{ecm\_params.mat} and \texttt{drivecycles\_meas.mat} into \texttt{data/}.
    \item Verify workspace variables exist:
    \begin{verbatim}
    load('data/ecm_params.mat')
    whos
    \end{verbatim}
    You should see: \texttt{SoC\_bp\_ocv} $(101\!\times\!1)$, \texttt{SoC\_bp\_ecm} $(21\!\times\!1)$, \texttt{T\_bp} $(1\!\times\!7)$, \texttt{OCV\_data}, \texttt{R0\_data}, etc.
    \item Open your existing model to review it:
    \begin{verbatim}
    open_system('models/cell_ecm_2rc')
    \end{verbatim}
    Take note of the block structure: inputs (From Workspace / Inports), OCV LUT, R0 LUT, R1/C1/R2/C2 LUTs, Coulomb counting integrator, output ports.
\end{enumerate}
\end{guistep}

%###############################################################################
%  PHASE 1 — MODULAR PACK ARCHITECTURE
%###############################################################################
\newpage
\section{Modular Pack Architecture}

%───────────────────────────────────────────────────────────────────────────────
\subsection{Create the Cell Parameter Configuration Script}
%───────────────────────────────────────────────────────────────────────────────

This script generates the 96-cell parameter table that will be used by all cell instances.

\begin{scriptbox}{Create \texttt{scripts/build\_cell\_params.m}}
In MATLAB: \textbf{Home} $\to$ \textbf{New Script}. Paste the following, save as \texttt{scripts/build\_cell\_params.m}:

\begin{verbatim}
% build_cell_params.m
% Generates 96-cell parameter table and saves to data/cell_params.mat

scriptDir = fileparts(mfilename('fullpath'));
dataDir   = fullfile(scriptDir, '..', 'data');

N_cells = 96;
N_modules = 8;
cells_per_module = 12;

% --- Cell types ---
%  Type 1: Nominal  (70%)  Q=60Ah, R0_scale=1.00, dSoC=0%
%  Type 2: Weak     (15%)  Q=57Ah, R0_scale=1.10, dSoC=-1.5%
%  Type 3: Strong   (15%)  Q=63Ah, R0_scale=0.92, dSoC=+1.0%

rng(42);  % reproducible

type_id = ones(N_cells, 1);           % default: nominal
idx_shuffle = randperm(N_cells);
type_id(idx_shuffle(1:14))  = 2;      % 14 weak
type_id(idx_shuffle(15:28)) = 3;      % 14 strong

Q_Ah      = zeros(N_cells, 1);
R0_scale  = zeros(N_cells, 1);
dSoC_pct  = zeros(N_cells, 1);

for i = 1:N_cells
    switch type_id(i)
        case 1, Q_Ah(i)=60; R0_scale(i)=1.00; dSoC_pct(i)= 0.0;
        case 2, Q_Ah(i)=57; R0_scale(i)=1.10; dSoC_pct(i)=-1.5;
        case 3, Q_Ah(i)=63; R0_scale(i)=0.92; dSoC_pct(i)=+1.0;
    end
end

% Build per-module cell parameter arrays (8 structs of 12 cells each)
module_id = reshape(1:N_cells, cells_per_module, N_modules);

cell_params = table(type_id, Q_Ah, R0_scale, dSoC_pct, ...
    'VariableNames', {'type','Q_Ah','R0_scale','dSoC_pct'});

% Save as matrices for Simulink (each row = 1 cell)
cell_Q_Ah     = Q_Ah;
cell_R0_scale = R0_scale;
cell_dSoC_pct = dSoC_pct;

save(fullfile(dataDir, 'cell_params.mat'), ...
    'cell_params', 'cell_Q_Ah', 'cell_R0_scale', 'cell_dSoC_pct', ...
    'N_cells', 'N_modules', 'cells_per_module', 'module_id');

fprintf('cell_params.mat saved (%d cells: %d nom, %d weak, %d strong)\n', ...
    N_cells, sum(type_id==1), sum(type_id==2), sum(type_id==3));
\end{verbatim}
Run it: \texttt{run('scripts/build\_cell\_params.m')}
\end{scriptbox}

%───────────────────────────────────────────────────────────────────────────────
\subsection{Create the Parameterised Cell Subsystem}
%───────────────────────────────────────────────────────────────────────────────

We will create a new model \texttt{cell\_ecm\_thermal.slx} based on the existing \texttt{cell\_ecm\_2rc.slx}, adding mask parameters and a thermal output.

\begin{guistep}{1.2.1 — Duplicate the existing cell model}
\begin{enumerate}
    \item In MATLAB Command Window:
    \begin{verbatim}
    open_system('models/cell_ecm_2rc')
    \end{verbatim}
    \item \textbf{File} $\to$ \textbf{Save As} $\to$ save as \texttt{models/cell\_ecm\_thermal.slx}.
    \item You now have a copy to modify. Keep it open.
\end{enumerate}
\end{guistep}

\begin{guistep}{1.2.2 — Convert the top level into a reusable Subsystem}
\begin{enumerate}
    \item Select \textbf{all blocks} in the canvas: \texttt{Ctrl+A}.
    \item Right-click $\to$ \textbf{Create Subsystem from Selection}.
    \item Double-click the new subsystem, rename it to \texttt{CellECM} (click on the label text).
    \item Inside the subsystem, verify you see the original ECM blocks plus Inport/Outport blocks that were auto-created.
\end{enumerate}
\end{guistep}

\begin{guistep}{1.2.3 — Set up Inports (inside CellECM subsystem)}
Verify or create these three Inport blocks at the left side of the subsystem canvas:
\begin{enumerate}
    \item \textbf{Inport 1}: Name = \texttt{I\_cell} \quad (cell current [A])
    \item \textbf{Inport 2}: Name = \texttt{T\_cell} \quad (cell temperature [\si{\degreeCelsius}])
    \item \textbf{Inport 3}: Name = \texttt{SoC\_init} \quad (initial SoC [\%])
\end{enumerate}
To rename: double-click the inport label text, type the new name.\\
If inports already exist from the original model, just rename them to match.
\end{guistep}

\begin{guistep}{1.2.4 — Set up Outports (inside CellECM subsystem)}
At the right side, create/verify these Outport blocks:
\begin{enumerate}
    \item \textbf{Outport 1}: Name = \texttt{V\_cell} \quad (cell voltage [V])
    \item \textbf{Outport 2}: Name = \texttt{SoC} \quad (state of charge [\%])
    \item \textbf{Outport 3}: Name = \texttt{Q\_gen} \quad (heat generation [W]) — \textbf{NEW}
\end{enumerate}
For the new \texttt{Q\_gen} output:
\begin{enumerate}
    \item From the Library Browser (\texttt{Ctrl+Shift+L}), drag a \textbf{Simulink $\to$ Sinks $\to$ Out1} block onto the canvas.
    \item Rename it to \texttt{Q\_gen}.
    \item To compute $Q_\text{gen} = I^2 \times R_0$:
    \begin{itemize}
        \item Add a \textbf{Math $\to$ Math Function} block, set function to \texttt{square}. Connect \texttt{I\_cell} to its input.
        \item Add a \textbf{Math $\to$ Product} block. Connect the squared current and the R0 LUT output to its two inputs.
        \item Connect the Product output to the \texttt{Q\_gen} outport.
    \end{itemize}
\end{enumerate}
\end{guistep}

\begin{guistep}{1.2.5 — Add Mask Parameters to CellECM}
Navigate back to the top level (click the model name in the breadcrumb bar).
\begin{enumerate}
    \item Right-click the \texttt{CellECM} subsystem block $\to$ \textbf{Mask} $\to$ \textbf{Create Mask...}
    \item Go to the \textbf{Parameters \& Dialog} tab.
    \item Click \textbf{Edit} (pencil icon). Add three parameters by clicking the ``+'' button:

    \begin{center}
    \begin{tabular}{llll}
        \toprule
        \textbf{Prompt} & \textbf{Variable} & \textbf{Type} & \textbf{Default} \\
        \midrule
        Capacity [Ah]      & \texttt{Q\_Ah}      & \texttt{edit} & \texttt{60} \\
        R0 scale factor     & \texttt{R0\_scale}  & \texttt{edit} & \texttt{1.0} \\
        SoC offset [\%]     & \texttt{dSoC\_pct}  & \texttt{edit} & \texttt{0} \\
        \bottomrule
    \end{tabular}
    \end{center}

    \item Click \textbf{OK} to close the mask editor.
\end{enumerate}
\end{guistep}

\begin{guistep}{1.2.6 — Wire mask parameters inside the subsystem}
Double-click \texttt{CellECM} to enter it. The mask variables (\texttt{Q\_Ah}, \texttt{R0\_scale}, \texttt{dSoC\_pct}) are now accessible inside.
\begin{enumerate}
    \item \textbf{Coulomb counting integrator:}
    \begin{itemize}
        \item Find the integrator that computes SoC. It should use the nominal capacity $Q$ to convert current to SoC.
        \item Locate the Gain block (or Constant) that holds \texttt{Q\_nom\_Ah} (currently = 60).
        \item Double-click it, change its value from \texttt{Q\_nom\_Ah} to \texttt{Q\_Ah}.
    \end{itemize}

    \item \textbf{Initial SoC:}
    \begin{itemize}
        \item Find the integrator's Initial Condition (IC) source. It uses \texttt{SoC\_init} from the inport.
        \item Add a \textbf{Math $\to$ Add} block between the \texttt{SoC\_init} inport and the integrator IC.
        \item Add a \textbf{Sources $\to$ Constant} block, set value = \texttt{dSoC\_pct}.
        \item Connect: \texttt{SoC\_init} + \texttt{dSoC\_pct} $\to$ integrator IC.
    \end{itemize}

    \item \textbf{R0 scaling:}
    \begin{itemize}
        \item Find the R0 lookup table output signal.
        \item Add a \textbf{Math $\to$ Product} block right after the R0 LUT output.
        \item Add a \textbf{Sources $\to$ Constant} block, set value = \texttt{R0\_scale}.
        \item Connect: $R_{0,\text{LUT}} \times \texttt{R0\_scale}$ $\to$ continues to V calculation.
    \end{itemize}
\end{enumerate}
\end{guistep}

\begin{guistep}{1.2.7 — Verify the cell model}
\begin{enumerate}
    \item Press \textbf{Ctrl+D} to update the diagram (checks for wiring errors).
    \item Add temporary test signals at top level:
    \begin{itemize}
        \item Add a \textbf{Signal Builder} or \textbf{From Workspace} block for \texttt{I\_cell} (constant -20A).
        \item Add a \textbf{Constant} for \texttt{T\_cell} = 25.
        \item Add a \textbf{Constant} for \texttt{SoC\_init} = 80.
    \end{itemize}
    \item Set solver: \textbf{Modeling} $\to$ \textbf{Model Settings} $\to$ \textbf{Solver}:
    \begin{itemize}
        \item Type = Fixed-step, Solver = ode4, Fixed-step size = 0.1
        \item Stop time = 100
    \end{itemize}
    \item Click \textbf{Run} (green play button). Check Scope blocks for reasonable V\_cell ($\sim$3.5--4.1 V) and SoC (decreasing).
    \item Test with mask: double-click \texttt{CellECM} block $\to$ change \texttt{Q\_Ah} to 57, re-run $\to$ SoC should drop faster.
    \item \textbf{Save}: \texttt{Ctrl+S}.
\end{enumerate}
\end{guistep}

%───────────────────────────────────────────────────────────────────────────────
\subsection{Build the Module Subsystem (12 cells in series)}
%───────────────────────────────────────────────────────────────────────────────

\begin{guistep}{1.3.1 — Create new model file}
\begin{enumerate}
    \item \textbf{File} $\to$ \textbf{New} $\to$ \textbf{Blank Model}.
    \item \textbf{File} $\to$ \textbf{Save As} $\to$ \texttt{models/module\_12s.slx}.
\end{enumerate}
\end{guistep}

\begin{guistep}{1.3.2 — Add Inports}
Drag from Library Browser (\textbf{Simulink $\to$ Sources $\to$ In1}):
\begin{enumerate}
    \item \texttt{I\_module} — module current [A] (same as cell current, series connection)
    \item \texttt{T\_cells} — 12$\times$1 vector of cell temperatures [\si{\degreeCelsius}]
    \item \texttt{SoC\_init} — scalar initial SoC [\%] (base value, offsets applied per cell)
    \item \texttt{cell\_Q} — 12$\times$1 vector of cell capacities [Ah]
    \item \texttt{cell\_R0s} — 12$\times$1 vector of R0 scale factors
    \item \texttt{cell\_dSoC} — 12$\times$1 vector of SoC offsets [\%]
\end{enumerate}
Arrange them vertically on the left side of the canvas.
\end{guistep}

\begin{guistep}{1.3.3 — Add Outports}
Drag \textbf{Simulink $\to$ Sinks $\to$ Out1}:
\begin{enumerate}
    \item \texttt{V\_module} — scalar module voltage [V]
    \item \texttt{V\_cells} — 12$\times$1 cell voltage vector [V]
    \item \texttt{SoC\_cells} — 12$\times$1 cell SoC vector [\%]
    \item \texttt{Q\_gen\_cells} — 12$\times$1 heat generation vector [W]
\end{enumerate}
Arrange on the right side.
\end{guistep}

\begin{guistep}{1.3.4 — Place 12 cell instances using For Each Subsystem}
Instead of manually copying 12 cells (error-prone), use a \textbf{For Each Subsystem}:
\begin{enumerate}
    \item From Library Browser: \textbf{Simulink $\to$ Ports \& Subsystems $\to$ For Each Subsystem}. Drag onto canvas.
    \item Double-click to enter the For Each Subsystem.
    \item Inside, you see \textbf{For Each} block at top. Double-click it:
    \begin{itemize}
        \item Set \textbf{Partition Dimension} = \texttt{1} and \textbf{Partition Width} = \texttt{1} for each input that is a vector.
    \end{itemize}
    \item Inside this subsystem, place a \textbf{Model Reference} or copy the \texttt{CellECM} subsystem:
    \begin{itemize}
        \item \textbf{Option A (Model Reference — recommended):}\\
              Library Browser $\to$ \textbf{Simulink $\to$ Ports \& Subsystems $\to$ Model}. Drag it in.\\
              Double-click $\to$ set \textbf{Model name} = \texttt{cell\_ecm\_thermal}.\\
              This references \texttt{cell\_ecm\_thermal.slx} without duplicating it.
        \item \textbf{Option B (Copy subsystem):}\\
              Open \texttt{cell\_ecm\_thermal.slx}, copy the CellECM subsystem block (\texttt{Ctrl+C}), paste here (\texttt{Ctrl+V}).
    \end{itemize}
    \item Wire the For Each inports to the cell block inputs:
    \begin{itemize}
        \item \texttt{In1} (I\_cell) $\leftarrow$ the module current (scalar, same for all iterations)
        \item \texttt{In2} (T\_cell) $\leftarrow$ partitioned from T\_cells vector (one element per iteration)
        \item \texttt{In3} (SoC\_init) $\leftarrow$ scalar, same for all
        \item Pass \texttt{cell\_Q}, \texttt{cell\_R0s}, \texttt{cell\_dSoC} as partitioned inputs
    \end{itemize}
    \item Wire cell outputs to For Each outports:
    \begin{itemize}
        \item \texttt{Out1} = V\_cell (will be concatenated to 12$\times$1)
        \item \texttt{Out2} = SoC (concatenated to 12$\times$1)
        \item \texttt{Out3} = Q\_gen (concatenated to 12$\times$1)
    \end{itemize}
    \item In the \textbf{For Each} block settings:
    \begin{itemize}
        \item For \texttt{I\_module} and \texttt{SoC\_init}: uncheck ``Partition input'' (broadcast scalar).
        \item For \texttt{T\_cells}, \texttt{cell\_Q}, \texttt{cell\_R0s}, \texttt{cell\_dSoC}: check ``Partition input'', Dimension=1, Width=1.
    \end{itemize}
\end{enumerate}
\end{guistep}

\begin{guistep}{1.3.5 — Compute module voltage}
Back at the module top level:
\begin{enumerate}
    \item The V\_cells output from the For Each block is a 12$\times$1 vector.
    \item Add a \textbf{Math $\to$ Sum} block (or \textbf{Sum of Elements}).
    \begin{itemize}
        \item Double-click $\to$ set \textbf{Function} = \texttt{Sum of Elements} (or ``List of signs'' = \texttt{+}).
        \item Check \textbf{Over} = \texttt{All dimensions} (sums the full vector).
    \end{itemize}
    \item Connect V\_cells $\to$ Sum of Elements $\to$ \texttt{V\_module} outport.
    \item Also branch the V\_cells signal directly to the \texttt{V\_cells} outport (right-click wire $\to$ branch).
    \item Connect SoC\_cells and Q\_gen\_cells outports similarly.
    \item \textbf{Ctrl+D} to update diagram. \textbf{Ctrl+S} to save.
\end{enumerate}
\end{guistep}

\begin{warnbox}
If using Model Reference for the cell, ensure \texttt{cell\_ecm\_thermal.slx} is on the MATLAB path. Add it: \texttt{addpath('models/')}.
\end{warnbox}

%───────────────────────────────────────────────────────────────────────────────
\subsection{Build the Pack Top-Level (8 modules)}
%───────────────────────────────────────────────────────────────────────────────

\begin{guistep}{1.4.1 — Create new model}
\begin{enumerate}
    \item \textbf{File} $\to$ \textbf{New} $\to$ \textbf{Blank Model}.
    \item Save as \texttt{models/pack\_96s.slx}.
\end{enumerate}
\end{guistep}

\begin{guistep}{1.4.2 — Pack Inports}
\begin{enumerate}
    \item \texttt{I\_pack} — pack current [A]
    \item \texttt{T\_amb} — ambient temperature [\si{\degreeCelsius}]
    \item \texttt{SoC\_init} — scalar initial SoC [\%]
    \item \texttt{T\_coolant} — coolant temperature [\si{\degreeCelsius}] (from thermal model, to be connected later)
\end{enumerate}
\end{guistep}

\begin{guistep}{1.4.3 — Pack Outports}
\begin{enumerate}
    \item \texttt{V\_pack} — pack voltage [V]
    \item \texttt{V\_cells\_96} — 96$\times$1 all cell voltages [V]
    \item \texttt{SoC\_cells\_96} — 96$\times$1 all cell SoCs [\%]
    \item \texttt{T\_cells\_96} — 96$\times$1 all cell temperatures [\si{\degreeCelsius}]
    \item \texttt{Q\_gen\_96} — 96$\times$1 all cell heat generation [W]
\end{enumerate}
\end{guistep}

\begin{guistep}{1.4.4 — Use a second For Each Subsystem for 8 modules}
\begin{enumerate}
    \item Drag a \textbf{For Each Subsystem} onto the pack canvas.
    \item Inside, place a Model Reference to \texttt{module\_12s} (or copy the module subsystem).
    \item Configure the For Each block:
    \begin{itemize}
        \item \texttt{I\_pack}: broadcast (no partition) — same current for all modules.
        \item \texttt{SoC\_init}: broadcast.
        \item \textbf{Cell parameters}: reshaped as 12$\times$8 matrices at pack level. Each iteration gets one column.
    \end{itemize}
    \item For the cell parameter routing (before the For Each block):
    \begin{itemize}
        \item Add a \textbf{Constant} block, set value = \texttt{reshape(cell\_Q\_Ah, 12, 8)} — this creates a $12\!\times\!8$ matrix.
        \item Partition along dimension 2, width 1 $\to$ each iteration gets 12$\times$1.
        \item Repeat for \texttt{cell\_R0\_scale} and \texttt{cell\_dSoC\_pct}.
    \end{itemize}
    \item Outputs: V\_module (8$\times$1 scalars), V\_cells (12$\times$8), SoC\_cells (12$\times$8), Q\_gen (12$\times$8).
\end{enumerate}
\end{guistep}

\begin{guistep}{1.4.5 — Compute pack voltage and flatten sensor vectors}
\begin{enumerate}
    \item \textbf{V\_pack}: Add a \textbf{Sum of Elements} block $\to$ sum the 8 module voltages.
    \item \textbf{Flatten 12$\times$8 to 96$\times$1}: For each matrix output (V\_cells, SoC\_cells, Q\_gen):
    \begin{itemize}
        \item Add a \textbf{Signal Processing $\to$ Reshape} block (or \textbf{Math $\to$ Reshape}).
        \item Set output dimensions = \texttt{[96, 1]}.
    \end{itemize}
    \item Connect to the corresponding outports.
    \item \textbf{Ctrl+D}, \textbf{Ctrl+S}.
\end{enumerate}
\end{guistep}

\begin{guistep}{1.4.6 — Quick sanity test of the pack}
\begin{enumerate}
    \item At the pack top level, add temporary test inputs:
    \begin{itemize}
        \item \texttt{I\_pack} = Constant -30 (30A discharge)
        \item \texttt{T\_amb} = Constant 25
        \item \texttt{SoC\_init} = Constant 80
        \item \texttt{T\_coolant} = Constant 25
    \end{itemize}
    \item Before running, load parameters into the workspace:
    \begin{verbatim}
    load('data/ecm_params.mat')
    load('data/cell_params.mat')
    \end{verbatim}
    \item Add \textbf{Scope} blocks on V\_pack and SoC\_cells\_96. Run for 100~s.
    \item Expected: V\_pack $\approx 96 \times 3.9 \approx \SI{374}{V}$. SoC decreasing. Weak cells diverge slightly.
    \item Delete temporary test blocks after verification.
\end{enumerate}
\end{guistep}


%###############################################################################
%  PHASE 2 — THERMAL MODEL
%###############################################################################
\newpage
\section{Thermal Model}

%───────────────────────────────────────────────────────────────────────────────
\subsection{Add thermal node inside CellECM subsystem}
%───────────────────────────────────────────────────────────────────────────────

\begin{guistep}{2.1.1 — Open the cell model and add thermal inports}
\begin{enumerate}
    \item Open \texttt{cell\_ecm\_thermal.slx}. Double-click into the \texttt{CellECM} subsystem.
    \item Add a new Inport: \texttt{T\_coolant} (port 4).
    \item Add a new Inport: \texttt{T\_amb} (port 5).
\end{enumerate}
\end{guistep}

\begin{guistep}{2.1.2 — Build the thermal ODE subsystem}
Inside the CellECM subsystem, create a new subsystem for the thermal equation:
\begin{enumerate}
    \item Right-click canvas $\to$ \textbf{Add Block} $\to$ \textbf{Subsystem}. Name it \texttt{ThermalNode}.
    \item Double-click to enter \texttt{ThermalNode}.
    \item Create these Inports:
    \begin{itemize}
        \item \texttt{Q\_gen} (heat generation from $I^2 R_0$) [W]
        \item \texttt{T\_coolant} [\si{\degreeCelsius}]
        \item \texttt{T\_amb} [\si{\degreeCelsius}]
        \item \texttt{T\_cell\_init} [\si{\degreeCelsius}] (for integrator IC)
    \end{itemize}
    \item Create one Outport: \texttt{T\_cell} [\si{\degreeCelsius}]
    \item \textbf{Build the thermal equation} ($\dot{T} = \frac{1}{m c_p}[Q_\text{gen} - h_\text{cool}(T - T_\text{cool}) - h_\text{amb}(T - T_\text{amb})]$):

    \textit{Step by step block placement:}

    \begin{enumerate}[label=\alph*)]
        \item Add \textbf{Integrator} block (Continuous $\to$ Integrator):
        \begin{itemize}
            \item Double-click $\to$ Initial condition source = \texttt{external}.
            \item Connect \texttt{T\_cell\_init} to the IC input (bottom port).
            \item Output = $T_\text{cell}$.
        \end{itemize}

        \item Add \textbf{Subtract} block \#1 (Math $\to$ Add, signs = \texttt{+-}):
        \begin{itemize}
            \item Inputs: $T_\text{cell}$ (from integrator output, feed back) and $T_\text{coolant}$.
            \item Output: $(T_\text{cell} - T_\text{coolant})$.
        \end{itemize}

        \item Add \textbf{Gain} block \#1: value = \texttt{h\_cool} (we will define this as a workspace variable).
        \begin{itemize}
            \item Connect input from Subtract \#1.
            \item Output: $h_\text{cool} \cdot (T_\text{cell} - T_\text{coolant})$.
        \end{itemize}

        \item Add \textbf{Subtract} block \#2: Inputs: $T_\text{cell}$, $T_\text{amb}$. Output: $(T_\text{cell} - T_\text{amb})$.

        \item Add \textbf{Gain} block \#2: value = \texttt{h\_amb}.

        \item Add a \textbf{Sum} block with signs = \texttt{+--}:
        \begin{itemize}
            \item Input 1 (+): \texttt{Q\_gen}
            \item Input 2 (--): output of Gain \#1 ($h_\text{cool}$ term)
            \item Input 3 (--): output of Gain \#2 ($h_\text{amb}$ term)
        \end{itemize}

        \item Add \textbf{Gain} block \#3: value = \texttt{1/(m\_cell * cp\_cell)}
        \begin{itemize}
            \item Connect from Sum output $\to$ Gain \#3 $\to$ Integrator input.
        \end{itemize}

        \item Connect Integrator output $\to$ \texttt{T\_cell} outport.

        \item \textbf{Feedback loop}: Branch the integrator output back to Subtract \#1 and Subtract \#2 positive inputs.
    \end{enumerate}
\end{enumerate}
\end{guistep}

\begin{guistep}{2.1.3 — Wire ThermalNode into CellECM}
Back in the CellECM subsystem:
\begin{enumerate}
    \item Connect \texttt{Q\_gen} signal (computed as $I^2 \cdot R_0$) $\to$ ThermalNode/Q\_gen.
    \item Connect \texttt{T\_coolant} inport $\to$ ThermalNode/T\_coolant.
    \item Connect \texttt{T\_amb} inport $\to$ ThermalNode/T\_amb.
    \item For \texttt{T\_cell\_init}: use the first sample of \texttt{T\_cell} input (add a \textbf{Memory} or \textbf{IC} block set to \texttt{T\_cell} inport at $t=0$).
    \item \textbf{Critical change}: The ECM LUTs (OCV, R0, R1, C1, R2, C2) currently use \texttt{T\_cell} from the inport. Replace: disconnect them from Inport and connect to the \textbf{ThermalNode output} \texttt{T\_cell} instead. This closes the electro-thermal loop.
    \item Rename the old \texttt{T\_cell} inport to \texttt{T\_cell\_init} (used only for initial condition now).
    \item Add a new outport: \texttt{T\_cell\_out} connected to the ThermalNode output (for monitoring).
\end{enumerate}
\end{guistep}

\begin{guistep}{2.1.4 — Update cell mask}
Right-click \texttt{CellECM} $\to$ \textbf{Mask} $\to$ \textbf{Edit Mask}. Add one parameter:
\begin{center}
\begin{tabular}{llll}
    \toprule
    \textbf{Prompt} & \textbf{Variable} & \textbf{Type} & \textbf{Default} \\
    \midrule
    Initial cell temp [\si{\degreeCelsius}] & \texttt{T\_cell\_init\_val} & edit & \texttt{25} \\
    \bottomrule
\end{tabular}
\end{center}
Click \textbf{OK}. Save model.
\end{guistep}

%───────────────────────────────────────────────────────────────────────────────
\subsection{Build the Coolant Loop Subsystem}
%───────────────────────────────────────────────────────────────────────────────

This is a standalone subsystem at the pack level.

\begin{guistep}{2.2.1 — Create subsystem in pack\_96s.slx}
\begin{enumerate}
    \item Open \texttt{pack\_96s.slx}.
    \item Right-click canvas $\to$ \textbf{Add Block} $\to$ \textbf{Subsystem}. Name: \texttt{CoolantLoop}.
    \item Double-click to enter.
\end{enumerate}
\end{guistep}

\begin{guistep}{2.2.2 — CoolantLoop I/O}
\textbf{Inports:}
\begin{enumerate}
    \item \texttt{Q\_cells\_total} — scalar: total cell heat [W] (sum of all 96 Q\_gen)
    \item \texttt{Q\_heater} — heater power [W]
    \item \texttt{Q\_chiller} — chiller cooling power [W]
    \item \texttt{T\_amb} — ambient temperature [\si{\degreeCelsius}]
    \item \texttt{T\_cool\_init} — initial coolant temperature [\si{\degreeCelsius}]
\end{enumerate}
\textbf{Outport:}
\begin{enumerate}
    \item \texttt{T\_coolant} — coolant temperature [\si{\degreeCelsius}]
\end{enumerate}
\end{guistep}

\begin{guistep}{2.2.3 — Build coolant ODE}
Implements: $\dot{T}_\text{cool} = \frac{1}{m_c c_{p,c}} \left[ Q_\text{cells} + Q_\text{heater} - Q_\text{chiller} - h_\text{rad}(T_\text{cool} - T_\text{amb}) \right]$

\begin{enumerate}
    \item Add \textbf{Integrator}:
    \begin{itemize}
        \item Initial condition source = external. Connect \texttt{T\_cool\_init} to IC port.
        \item Output = $T_\text{coolant}$.
    \end{itemize}

    \item Add \textbf{Subtract}: $T_\text{coolant} - T_\text{amb}$.

    \item Add \textbf{Gain}: value = \texttt{h\_rad} (radiator conductance, workspace variable).

    \item Add \textbf{Sum} with signs = \texttt{++--}:
    \begin{itemize}
        \item (+) \texttt{Q\_cells\_total}
        \item (+) \texttt{Q\_heater}
        \item (--) \texttt{Q\_chiller}
        \item (--) $h_\text{rad}(T_\text{cool} - T_\text{amb})$
    \end{itemize}

    \item Add \textbf{Gain}: value = \texttt{1/(m\_coolant * cp\_coolant)}.

    \item Chain: Sum $\to$ Gain $\to$ Integrator $\to$ T\_coolant outport + feedback.

    \item \textbf{Ctrl+D}, \textbf{Ctrl+S}.
\end{enumerate}
\end{guistep}

%───────────────────────────────────────────────────────────────────────────────
\subsection{Create Thermal Parameters Script}
%───────────────────────────────────────────────────────────────────────────────

\begin{scriptbox}{Create \texttt{scripts/build\_thermal\_params.m}}
\begin{verbatim}
% build_thermal_params.m
% Defines thermal model parameters

scriptDir = fileparts(mfilename('fullpath'));
dataDir   = fullfile(scriptDir, '..', 'data');

% --- Cell thermal parameters ---
m_cell    = 0.8;      % cell mass [kg]
cp_cell   = 1000;     % specific heat [J/(kg·K)]
h_cool    = 10;       % cell-to-coolant conductance [W/K]
h_amb     = 1.0;      % cell-to-ambient conductance [W/K]

% --- Coolant loop parameters ---
m_coolant  = 5.0;     % coolant mass [kg]
cp_coolant = 3500;    % coolant specific heat [J/(kg·K)] (water-glycol)
h_rad      = 15;      % radiator-to-ambient conductance [W/K]

% --- Actuator limits ---
P_heater_max  = 5000;  % [W]
P_chiller_max = 3000;  % [W]
tau_heater    = 5;     % heater time constant [s]
tau_chiller   = 10;    % chiller time constant [s]

% --- Thermal manager thresholds ---
T_heat_on     = 5;     % start heating below this [°C]
T_heat_off    = 7;     % stop heating above this [°C] (hysteresis)
T_cool_on     = 35;    % start cooling above this [°C]
T_cool_off    = 33;    % stop cooling below this [°C] (hysteresis)
T_heat_target = 15;    % coolant target when heating [°C]
T_cool_target = 20;    % coolant target when cooling [°C]

save(fullfile(dataDir, 'thermal_params.mat'), ...
    'm_cell','cp_cell','h_cool','h_amb', ...
    'm_coolant','cp_coolant','h_rad', ...
    'P_heater_max','P_chiller_max','tau_heater','tau_chiller', ...
    'T_heat_on','T_heat_off','T_cool_on','T_cool_off', ...
    'T_heat_target','T_cool_target');

fprintf('thermal_params.mat saved\n');
\end{verbatim}
Run: \texttt{run('scripts/build\_thermal\_params.m')}
\end{scriptbox}


%###############################################################################
%  PHASE 3 — THERMAL MANAGEMENT SYSTEM
%###############################################################################
\newpage
\section{Thermal Management System (Plant Actuator)}

%───────────────────────────────────────────────────────────────────────────────
\subsection{Build the Heater subsystem}
%───────────────────────────────────────────────────────────────────────────────

\begin{guistep}{3.1.1 — Create Heater subsystem in pack\_96s.slx}
\begin{enumerate}
    \item In \texttt{pack\_96s.slx}, add a new Subsystem: \texttt{Heater}.
    \item \textbf{Inport}: \texttt{P\_cmd} — commanded heater power [W].
    \item \textbf{Outport}: \texttt{Q\_heater} — actual heater power [W].
    \item Inside, build:
    \begin{enumerate}[label=\alph*)]
        \item \textbf{Saturation} block: Lower = 0, Upper = \texttt{P\_heater\_max}.\\
              Connect \texttt{P\_cmd} $\to$ Saturation.
        \item \textbf{Transfer Function} block: Numerator = \texttt{[1]}, Denominator = \texttt{[tau\_heater, 1]}.\\
              (This gives $\frac{1}{\tau_h s + 1}$).
        \item Wire: \texttt{P\_cmd} $\to$ Saturation $\to$ Transfer Fcn $\to$ \texttt{Q\_heater}.
    \end{enumerate}
\end{enumerate}
\end{guistep}

%───────────────────────────────────────────────────────────────────────────────
\subsection{Build the Chiller subsystem}
%───────────────────────────────────────────────────────────────────────────────

\begin{guistep}{3.2.1 — Create Chiller subsystem}
\begin{enumerate}
    \item Same as the Heater, but:
    \begin{itemize}
        \item Saturation Upper = \texttt{P\_chiller\_max}.
        \item Transfer Fcn denominator = \texttt{[tau\_chiller, 1]}.
    \end{itemize}
    \item \textbf{Inport}: \texttt{P\_cmd}. \textbf{Outport}: \texttt{Q\_chiller}.
\end{enumerate}
\end{guistep}

%───────────────────────────────────────────────────────────────────────────────
\subsection{Build the Thermal Manager Stub Controller}
%───────────────────────────────────────────────────────────────────────────────

\begin{guistep}{3.3.1 — Create ThermalManagerStub subsystem}
\begin{enumerate}
    \item Add subsystem \texttt{ThermalManagerStub} in \texttt{pack\_96s.slx}.
    \item \textbf{Inports}:
    \begin{itemize}
        \item \texttt{T\_cell\_min} — minimum cell temperature [\si{\degreeCelsius}]
        \item \texttt{T\_cell\_max} — maximum cell temperature [\si{\degreeCelsius}]
        \item \texttt{T\_coolant} — current coolant temperature [\si{\degreeCelsius}]
    \end{itemize}
    \item \textbf{Outports}:
    \begin{itemize}
        \item \texttt{P\_heat\_cmd} — heater command [W]
        \item \texttt{P\_chill\_cmd} — chiller command [W]
    \end{itemize}
\end{enumerate}
\end{guistep}

\begin{guistep}{3.3.2 — Implement heating logic with hysteresis}
Inside \texttt{ThermalManagerStub}:

\textbf{Heating path:}
\begin{enumerate}
    \item Add a \textbf{Relay} block (Discontinuities $\to$ Relay):
    \begin{itemize}
        \item Switch on point = \texttt{T\_heat\_on} (5\si{\degreeCelsius})
        \item Switch off point = \texttt{T\_heat\_off} (7\si{\degreeCelsius})
        \item Output when on = 1, Output when off = 0
        \item \textbf{Direction}: Set to ``switch on when input falls below on-point''.\\
              (Since we want heating when $T < 5$\si{\degreeCelsius}, negate the input first.)
    \end{itemize}
    \item Add a \textbf{Gain} block with \texttt{-1} before the Relay input (since Relay triggers on rising edge by default, we negate so cold $\to$ high).
    \item The Relay outputs 0 or 1. Multiply by a \textbf{P controller}:
    \begin{itemize}
        \item Add a \textbf{Subtract}: $T_\text{heat\_target} - T_\text{coolant}$.
        \item Add a \textbf{Gain}: value = \texttt{P\_heater\_max / 10} (proportional gain, approximate).
        \item Add a \textbf{Saturation}: [0, \texttt{P\_heater\_max}].
        \item Multiply Relay output $\times$ saturated P output $\to$ \texttt{P\_heat\_cmd}.
    \end{itemize}
    \item Alternatively, for simplicity: Relay output $\times$ \texttt{P\_heater\_max} $\to$ \texttt{P\_heat\_cmd} (bang-bang control).
\end{enumerate}

\textbf{Cooling path:} Same structure with:
\begin{itemize}
    \item Relay: on = \texttt{T\_cool\_on} (35), off = \texttt{T\_cool\_off} (33). No negation needed (triggers when hot).
    \item Output scaled by \texttt{P\_chiller\_max}.
\end{itemize}
\end{guistep}

\begin{guistep}{3.3.3 — Compute T\_cell\_min and T\_cell\_max}
At the pack level (outside the stub):
\begin{enumerate}
    \item Add \textbf{MinMax} block (set to ``min''), connect \texttt{T\_cells\_96} $\to$ MinMax $\to$ \texttt{T\_cell\_min}.
    \item Add another \textbf{MinMax} block (``max''), connect \texttt{T\_cells\_96} $\to$ MinMax $\to$ \texttt{T\_cell\_max}.
    \item Connect these to the \texttt{ThermalManagerStub} inports.
\end{enumerate}
\end{guistep}


%###############################################################################
%  PHASE 4 — INTEGRATION & VALIDATION
%###############################################################################
\newpage
\section{Integration \& Validation}

%───────────────────────────────────────────────────────────────────────────────
\subsection{Wire the complete feedback loop}
%───────────────────────────────────────────────────────────────────────────────

\begin{guistep}{4.1.1 — Connect all subsystems in pack\_96s.slx}
At the top level of \texttt{pack\_96s.slx}, wire the feedback loop:
\begin{enumerate}
    \item \texttt{Q\_gen\_96} (from For Each pack) $\to$ \textbf{Sum of Elements} $\to$ CoolantLoop/Q\_cells\_total.
    \item \texttt{T\_cells\_96} $\to$ MinMax blocks $\to$ ThermalManagerStub inputs.
    \item ThermalManagerStub/P\_heat\_cmd $\to$ Heater/P\_cmd.
    \item ThermalManagerStub/P\_chill\_cmd $\to$ Chiller/P\_cmd.
    \item Heater/Q\_heater $\to$ CoolantLoop/Q\_heater.
    \item Chiller/Q\_chiller $\to$ CoolantLoop/Q\_chiller.
    \item \texttt{T\_amb} inport $\to$ CoolantLoop/T\_amb.
    \item CoolantLoop/T\_coolant $\to$ \textbf{broadcast to all cells} via the pack inport \texttt{T\_coolant}.
    \item CoolantLoop/T\_coolant $\to$ ThermalManagerStub/T\_coolant.
    \item \textbf{Close the loop}: Delete the \texttt{T\_coolant} top-level inport. The coolant temperature is now internally computed.
\end{enumerate}

Final top-level inports should be:
\begin{itemize}
    \item \texttt{I\_pack}(t), \texttt{T\_amb}(t), \texttt{SoC\_init}
\end{itemize}

\textbf{Ctrl+D} to verify diagram. Fix any dimension mismatches or algebraic loops.
\end{guistep}

\begin{warnbox}
The thermal feedback creates a loop: Cells $\to$ Q\_gen $\to$ Coolant $\to$ T\_coolant $\to$ Cells. This is not an algebraic loop because both the cell integrator and coolant integrator introduce state delays. Simulink should handle it automatically. If you get an algebraic loop error, add a \textbf{Memory} block on the \texttt{T\_coolant} signal returning to cells.
\end{warnbox}

%───────────────────────────────────────────────────────────────────────────────
\subsection{Create the pack validation harness script}
%───────────────────────────────────────────────────────────────────────────────

\begin{scriptbox}{Create \texttt{scripts/validate\_pack.m}}
\begin{verbatim}
% validate_pack.m
% Validates pack_96s model against measurement trips

clear; clc;

scriptDir = fileparts(mfilename('fullpath'));
dataDir   = fullfile(scriptDir, '..', 'data');
plotDir   = fullfile(scriptDir, '..', 'plots', 'trips');
if ~exist(plotDir,'dir'), mkdir(plotDir); end

load(fullfile(dataDir, 'ecm_params.mat'));
load(fullfile(dataDir, 'cell_params.mat'));
load(fullfile(dataDir, 'thermal_params.mat'));
trips = load(fullfile(dataDir, 'drivecycles_meas.mat'));

TRIP_IDS = {'TripA01','TripA06','TripB01','TripB09','TripB14'};
mdl = 'pack_96s';
load_system(fullfile(scriptDir,'..','models',mdl));

results = table();

for ti = 1:numel(TRIP_IDS)
    tid  = TRIP_IDS{ti};
    trip = trips.(matlab.lang.makeValidName(tid));

    t_meas   = double(trip.time_s(:)) - double(trip.time_s(1));
    soc_init = double(trip.SoC_pct(1));

    % build input timeseries: [time, I_pack, T_amb]
    I_pack = double(trip.I_A(:));
    T_amb  = ones(size(t_meas)) * mean(double(trip.T_cell_C));

    simIn = Simulink.SimulationInput(mdl);
    simIn = simIn.setVariable('soc_init', soc_init);
    simIn = simIn.setModelParameter( ...
        'StopTime', num2str(t_meas(end)), ...
        'SolverType','Fixed-step','Solver','ode4', ...
        'FixedStep', num2str(dt));
    simIn = simIn.setExternalInput([t_meas, I_pack, T_amb]);

    simOut = sim(simIn);

    % Extract and compare
    V_pack_sim = simOut.yout(:,1);
    V_pack_ref = interp1(t_meas, double(trip.V_pack_V(:)), ...
                         simOut.tout, 'linear','extrap');

    err_mV = (V_pack_sim - V_pack_ref) * 1e3;

    row = {tid, rms(err_mV), mean(abs(err_mV)), ...
           max(abs(err_mV)), mean(err_mV)};
    results = [results; cell2table(row, ...
        'VariableNames',{'trip','RMSE_mV','MAE_mV', ...
                         'Max_mV','Bias_mV'})];

    fprintf('%s: RMSE=%.1f mV\n', tid, rms(err_mV));
end

close_system(mdl, 0);
disp(results);

writetable(results, fullfile(dataDir, ...
    'validation_results_pack.csv'));
\end{verbatim}
\end{scriptbox}

%───────────────────────────────────────────────────────────────────────────────
\subsection{Thermal parameter identification}
%───────────────────────────────────────────────────────────────────────────────

\begin{scriptbox}{Create \texttt{scripts/fit\_thermal\_params.m}}
\begin{verbatim}
% fit_thermal_params.m
% Fits h_cool, h_amb, m_cell*cp_cell using winter trips

scriptDir = fileparts(mfilename('fullpath'));
dataDir   = fullfile(scriptDir, '..', 'data');

load(fullfile(dataDir, 'ecm_params.mat'));
load(fullfile(dataDir, 'thermal_params.mat'));

% Load raw measurement data for thermal signals
T_raw = readtable(fullfile(scriptDir,'..','..','data','cleaned', ...
    'measurement','measurement_clean.csv'));

% Use 3-4 winter trips with large T variation
fit_trips = {'TripB09','TripB10','TripB14'};

% For each trip: extract T_cell, T_amb, I_cell, T_coolant
% Build a cost function: sum of squared (T_cell_sim - T_cell_meas)
% Optimise [h_cool, h_amb, m_cell*cp_cell] using fminsearch

% x0 = [h_cool, h_amb, mcp]
x0 = [10, 1.0, 800];
lb = [2, 0.1, 400];
ub = [30, 5, 2000];

% ... (implement cost function and optimisation here)
% After fitting, update thermal_params.mat

fprintf('Fit complete. Update thermal_params.mat with new values.\n');
\end{verbatim}
Skeleton — fill in the optimisation loop with your preferred method (\texttt{fminsearch}, \texttt{lsqnonlin}, or manual sweep).
\end{scriptbox}

%───────────────────────────────────────────────────────────────────────────────
\subsection{Run electrical validation}
%───────────────────────────────────────────────────────────────────────────────

\begin{guistep}{4.3.1 — Execute validation}
\begin{enumerate}
    \item In MATLAB Command Window:
    \begin{verbatim}
    load('data/ecm_params.mat')
    load('data/cell_params.mat')
    load('data/thermal_params.mat')
    run('scripts/validate_pack.m')
    \end{verbatim}
    \item Check the console output. Target: pack-level RMSE $< \SI{50}{mV}$ for warm trips.
    \item If RMSE is too high, debug:
    \begin{itemize}
        \item Run a single nominal cell and compare to previous single-cell validation.
        \item Check that R0\_scale and dSoC are applied correctly.
    \end{itemize}
\end{enumerate}
\end{guistep}

%───────────────────────────────────────────────────────────────────────────────
\subsection{Run thermal validation}
%───────────────────────────────────────────────────────────────────────────────

\begin{guistep}{4.4.1 — Compare modelled vs.\ measured temperature}
\begin{enumerate}
    \item After running \texttt{validate\_pack.m}, extract \texttt{T\_cells\_96} from simOut.
    \item Compute mean cell temperature: \texttt{T\_cell\_mean = mean(T\_cells\_96, 2)}.
    \item Plot against measured \texttt{battery\_temperature\_c}.
    \item Target: RMSE $< \SI{2}{\degreeCelsius}$.
    \item If thermal response is too slow/fast, adjust \texttt{h\_cool} and \texttt{m\_cell*cp\_cell} in \texttt{thermal\_params.mat}.
\end{enumerate}
\end{guistep}

%───────────────────────────────────────────────────────────────────────────────
\subsection{Verify cell imbalance}
%───────────────────────────────────────────────────────────────────────────────

\begin{guistep}{4.5.1 — Check SoC spread}
\begin{enumerate}
    \item Run a long trip (e.g. TripB14, $\sim$3800~s, large SoC swing 85\%$\to$35\%).
    \item Extract SoC\_cells\_96 from simOut (96 columns).
    \item Plot all 96 SoC traces vs.\ time:
    \begin{verbatim}
    plot(t_sim, SoC_cells_96)
    xlabel('Time [s]'); ylabel('SoC [%]');
    title('Cell SoC Spread');
    \end{verbatim}
    \item Expected: 3 distinct bands (nominal, weak, strong). Weak cells reach lower SoC first.
    \item SoC spread after 30+ min should be $\sim$2--5\%.
    \item Also check voltage spread:
    \begin{verbatim}
    plot(t_sim, V_cells_96)
    \end{verbatim}
    \item Weak cells (higher $R_0$) show larger voltage sag under load.
\end{enumerate}
\end{guistep}


%###############################################################################
%  PHASE 5 — BMS INTERFACE PREPARATION
%###############################################################################
\newpage
\section{BMS Interface Preparation}

%───────────────────────────────────────────────────────────────────────────────
\subsection{Define the Sensor Bus (Plant $\to$ BMS)}
%───────────────────────────────────────────────────────────────────────────────

\begin{guistep}{5.1.1 — Create a Bus Object for plant sensor outputs}
\begin{enumerate}
    \item In MATLAB Command Window, define the bus:
    \begin{verbatim}
    % Define sensor bus
    elems = Simulink.BusElement;
    elems(1) = Simulink.BusElement; elems(1).Name = 'V_cells';
        elems(1).Dimensions = [96 1];
    elems(2) = Simulink.BusElement; elems(2).Name = 'SoC_cells';
        elems(2).Dimensions = [96 1];
    elems(3) = Simulink.BusElement; elems(3).Name = 'T_cells';
        elems(3).Dimensions = [96 1];
    elems(4) = Simulink.BusElement; elems(4).Name = 'V_pack';
        elems(4).Dimensions = 1;
    elems(5) = Simulink.BusElement; elems(5).Name = 'I_pack';
        elems(5).Dimensions = 1;
    elems(6) = Simulink.BusElement; elems(6).Name = 'T_coolant';
        elems(6).Dimensions = 1;
    elems(7) = Simulink.BusElement; elems(7).Name = 'T_amb';
        elems(7).Dimensions = 1;
    SensorBus = Simulink.Bus;
    SensorBus.Elements = elems;
    \end{verbatim}
    \item In \texttt{pack\_96s.slx}, add a \textbf{Bus Creator} block.
    \begin{itemize}
        \item Set number of inputs = 7.
        \item Connect each signal to the bus creator.
        \item Double-click $\to$ check ``Specify properties via bus object'' $\to$ Bus name = \texttt{SensorBus}.
    \end{itemize}
    \item Connect Bus Creator output to a new outport: \texttt{SensorBus\_out}.
\end{enumerate}
\end{guistep}

%───────────────────────────────────────────────────────────────────────────────
\subsection{Define the Actuator Bus (BMS $\to$ Plant)}
%───────────────────────────────────────────────────────────────────────────────

\begin{guistep}{5.2.1 — Create actuator bus}
\begin{enumerate}
    \item In MATLAB:
    \begin{verbatim}
    elems2 = Simulink.BusElement;
    elems2(1) = Simulink.BusElement;
        elems2(1).Name = 'P_heat_cmd'; elems2(1).Dimensions = 1;
    elems2(2) = Simulink.BusElement;
        elems2(2).Name = 'P_chill_cmd'; elems2(2).Dimensions = 1;
    elems2(3) = Simulink.BusElement;
        elems2(3).Name = 'bal_switches';
        elems2(3).Dimensions = [96 1];
        elems2(3).DataType = 'boolean';
    ActuatorBus = Simulink.Bus;
    ActuatorBus.Elements = elems2;
    \end{verbatim}
    \item Add a new top-level Inport \texttt{ActuatorBus\_in} to \texttt{pack\_96s.slx}.
    \item Add a \textbf{Bus Selector} after it to extract: \texttt{P\_heat\_cmd}, \texttt{P\_chill\_cmd}, \texttt{bal\_switches}.
    \item Wire \texttt{P\_heat\_cmd} $\to$ Heater, \texttt{P\_chill\_cmd} $\to$ Chiller.
    \item \texttt{bal\_switches} will be used later (no effect yet — just expose the signal).
\end{enumerate}
\end{guistep}

%───────────────────────────────────────────────────────────────────────────────
\subsection{Build the BMS Stub}
%───────────────────────────────────────────────────────────────────────────────

\begin{guistep}{5.3.1 — Create BMS stub subsystem}
\begin{enumerate}
    \item In \texttt{pack\_96s.slx}, add a Subsystem: \texttt{BMS\_Stub}.
    \item \textbf{Inport}: \texttt{SensorBus} (receives the sensor bus from the pack).
    \item \textbf{Outport}: \texttt{ActuatorBus} (sends commands back to the pack).
    \item Inside:
    \begin{enumerate}[label=\alph*)]
        \item Add a \textbf{Bus Selector}: Extract \texttt{T\_cells} from SensorBus.
        \item Compute \textbf{min/max} of T\_cells (MinMax blocks).
        \item Extract \texttt{T\_coolant} from SensorBus.
        \item Copy the ThermalManagerStub logic inside (or reference it).
        \item Output: \texttt{P\_heat\_cmd}, \texttt{P\_chill\_cmd} from the thermal stub.
        \item For \texttt{bal\_switches}: add a \textbf{Constant} block = \texttt{false(96,1)}.
        \item Pack into an \textbf{ActuatorBus} using a Bus Creator.
    \end{enumerate}
    \item Wire: SensorBus\_out $\to$ BMS\_Stub/SensorBus $\to$ BMS\_Stub/ActuatorBus $\to$ ActuatorBus\_in.
    \item Now the original ThermalManagerStub at pack level can be replaced by the BMS\_Stub.
\end{enumerate}
\end{guistep}

\begin{guistep}{5.3.2 — Save bus definitions}
\begin{enumerate}
    \item In MATLAB:
    \begin{verbatim}
    save('data/bus_defs.mat', 'SensorBus', 'ActuatorBus');
    \end{verbatim}
    \item Add to the model initialization callback (\textbf{Modeling} $\to$ \textbf{Model Properties} $\to$ \textbf{Callbacks} $\to$ \textbf{InitFcn}):
    \begin{verbatim}
    load('data/ecm_params.mat');
    load('data/cell_params.mat');
    load('data/thermal_params.mat');
    load('data/bus_defs.mat');
    \end{verbatim}
    This ensures all parameters load automatically when the model opens.
\end{enumerate}
\end{guistep}


%###############################################################################
%  FINAL CHECK
%###############################################################################
\newpage
\section*{Final Checklist}
\addcontentsline{toc}{section}{Final Checklist}

\begin{center}
\renewcommand{\arraystretch}{1.3}
\begin{tabular}{p{0.4cm}p{12cm}c}
    \toprule
    & \textbf{Item} & \textbf{Done} \\
    \midrule
    1 & \texttt{scripts/build\_cell\_params.m} runs and saves \texttt{cell\_params.mat} & $\square$ \\
    2 & \texttt{cell\_ecm\_thermal.slx}: mask works (Q, R0\_scale, dSoC), Q\_gen output connected & $\square$ \\
    3 & \texttt{cell\_ecm\_thermal.slx}: ThermalNode subsystem computes $T_\text{cell}$ from $Q_\text{gen}$ & $\square$ \\
    4 & \texttt{module\_12s.slx}: For Each loop over 12 cells, V\_module = sum of V\_cells & $\square$ \\
    5 & \texttt{pack\_96s.slx}: 8 modules via For Each, V\_pack = sum, 96$\times$1 sensor vectors & $\square$ \\
    6 & CoolantLoop subsystem: single-node ODE, receives total Q\_gen + heater/chiller & $\square$ \\
    7 & Heater \& Chiller subsystems: saturation + first-order TF & $\square$ \\
    8 & ThermalManagerStub: Relay-based hysteresis for heat/cool modes & $\square$ \\
    9 & Full feedback loop wired and simulates without errors & $\square$ \\
    10 & \texttt{scripts/build\_thermal\_params.m} runs and saves \texttt{thermal\_params.mat} & $\square$ \\
    11 & Electrical validation: pack RMSE $<$ 50 mV (warm trips) & $\square$ \\
    12 & Thermal validation: $T_\text{cell}$ RMSE $<$ 2\si{\degreeCelsius} & $\square$ \\
    13 & Cell imbalance visible: SoC spread $\sim$2--5\% on long trips & $\square$ \\
    14 & SensorBus \& ActuatorBus defined and saved in \texttt{bus\_defs.mat} & $\square$ \\
    15 & BMS\_Stub block has correct I/O interface, balancing switches = OFF & $\square$ \\
    16 & Model InitFcn loads all .mat files automatically & $\square$ \\
    \bottomrule
\end{tabular}
\end{center}

\vspace{12pt}
\begin{center}
\fbox{\parbox{0.85\textwidth}{\centering
\textbf{After completing this guide}, the plant model is ready.\\
Next step: replace \texttt{BMS\_Stub} with the real BMS Simulink model\\
(SoC estimation, cell balancing, thermal control, fault detection).
}}
\end{center}

\end{document}
